\section{Introduction}
\label{sec:introduction}

\subsection{Purpose}
This Software Requirements Specification (SRS) establishes the complete, formal, and authoritative functional, behavioral, interface, and non-functional requirements for the \textbf{Midnight Compact Language Plugin for IntelliJ IDEA} (\texttt{dev.verloren.midnight}). 

This specification conforms strictly to the international standards \textbf{ISO/IEC/IEEE 29148:2018} (Systems and software engineering -- Life cycle processes -- Requirements engineering)~\cite{iso29148} and \textbf{IEEE Std 830-1998} (Recommended Practice for Software Requirements Specifications)~\cite{ieee830}. It serves as the baseline contract between software engineers, zero-knowledge (ZK) cryptographers, systems architects, quality assurance auditors, and the broader Midnight ecosystem development community.

\subsection{Document Conventions}
\label{sec:conventions}
This document enforces strict syntactic and semantic conventions to guarantee clarity, precision, and testability across all requirement definitions:

\begin{enumerate}
    \item \textbf{Normative Keyword Conventions (RFC 2119)}~\cite{rfc2119}:
    \begin{itemize}
        \item \textbf{SHALL}: Indicates an absolute, non-negotiable mandatory requirement of the system.
        \item \textbf{SHALL NOT}: Indicates an absolute prohibition.
        \item \textbf{SHOULD}: Indicates a recommended capability where valid reasons may exist in particular circumstances to ignore, after explicit architectural review.
        \item \textbf{MAY}: Indicates an optional or discretionary capability.
    \end{itemize}
    
    \item \textbf{Requirement Identification Scheme}:
    Every requirement throughout this document is assigned a persistent, globally unique identifier following the pattern:
    \begin{equation*}
        \mathbf{[REQ\text{-}SUBSYSTEM\text{-}CATEGORY\text{-}INDEX]}
    \end{equation*}
    where:
    \begin{itemize}
        \item \texttt{SUBSYSTEM} represents the architectural module: \texttt{LEX} (Lexer), \texttt{PRS} (Parser), \texttt{PSI} (Program Structure Interface), \texttt{RES} (Resolution), \texttt{TYP} (Type Engine), \texttt{IDE} (Navigation \& Completion), \texttt{INS} (Inspections), \texttt{CMP} (Compiler Tool Window), \texttt{MGR} (Toolchain Manager), \texttt{BAR} (Status Bar), \texttt{SEC} (Security), \texttt{PERF} (Performance).
        \item \texttt{CATEGORY} denotes the requirement type: \texttt{FUN} (Functional), \texttt{NFR} (Non-Functional), \texttt{INT} (Interface), or \texttt{CON} (Constraint).
        \item \texttt{INDEX} is a zero-padded three-digit sequential counter (e.g., \texttt{001}, \texttt{002}).
    \end{itemize}
    
    \item \textbf{EARS Syntax (Easy Approach to Requirements Syntax)}:
    All requirement statements follow standardized EARS patterns: Ubiquitous (``The system shall...''), Event-driven (``When <trigger>, the system shall...''), State-driven (``While <state>, the system shall...''), and Unwanted behavior (``If <error>, then the system shall...'').
\end{enumerate}

\subsection{Intended Audience and Reading Suggestions}
This document is targeted toward multiple distinct stakeholder groups:
\begin{itemize}
    \item \textbf{Plugin Core Developers}: Reference Section~\ref{sec:system_features}, Section~\ref{sec:external_interfaces}, and Section~\ref{sec:nonfunctional} for low-level AST contracts, threading requirements, and API constraints.
    \item \textbf{Smart Contract Developers \& ZK Engineers}: Focus on Section~\ref{sec:overall_description} and Section~\ref{sec:system_features} to understand IDE capabilities, semantic inspections, and compilation flows.
    \item \textbf{Security \& Compliance Auditors}: Review Section~\ref{sec:nonfunctional} (Security, Integrity, Isolation), and Section~\ref{sec:traceability_matrix} (Requirements Traceability).
    \item \textbf{Tooling \& DevOps Integrators}: Review Section~\ref{sec:external_interfaces} and Section~\ref{sec:appendices} for compiler version discovery, WSL integration, and CLI execution bindings.
\end{itemize}

\subsection{Product Scope and Strategic Objectives}
The \textbf{Midnight Compact Language Plugin} transforms JetBrains IntelliJ IDEA (and compatible IntelliJ Platform IDEs including IntelliJ IDEA Ultimate, Community Edition, and CLion) into an industrial-grade integrated development environment for the \textbf{Compact} smart contract language~\cite{compactSpec}. 

Compact is the domain-specific smart contract language developed for the \textbf{Midnight privacy-preserving blockchain}~\cite{midnightWhitepaper}. Midnight leverages zero-knowledge cryptography (ZK-SNARKs) to enable decentralized applications with selective data disclosure. Compact enables developers to author hybrid state machines that seamlessly interleave private witness operations (\texttt{witness}, \texttt{disclose}), public ledger state transitions (\texttt{ledger}, \texttt{sealed}), zero-knowledge circuits (\texttt{circuit}, \texttt{pure}), and domain cryptographic primitives (\texttt{Field}, \texttt{Uint}, \texttt{Cell}, \texttt{Map}, \texttt{Vector}).

The primary objectives of the plugin are:
\begin{enumerate}
    \item \textbf{Sub-millisecond Feedback Loop}: Provide instant lexical and syntactic error reporting, eliminating the latency of terminal compilation cycles.
    \item \textbf{Rigorous Static Semantics}: Detect illegal operations (e.g., mutating ledger state inside pure circuits, direct assignment of undisclosed witness state, or recursive circuit calls) directly within the editor via 10 static semantic inspections.
    \item \textbf{Deep Scoped Navigation}: Deliver fast, index-backed symbol resolution across local lexical scopes, module exports, type definitions, and multi-file project inclusions.
    \item \textbf{Isolated Multi-Version Compiler Lifecycle}: Manage multiple concurrent versions of the official Compact compiler toolchain in isolated user directories, providing automated background downloads, pragma SemVer evaluation, and a 1-click Remix-style build tool window.
\end{enumerate}

\begin{figure}[htbp]
    \centering
    % System Context Diagram (C4 Level 1)
\begin{tikzpicture}[
    scale=0.95, every node/.style={transform shape},
    actorNode/.style={
        rectangle, rounded corners=6pt,
        draw=srsDark, line width=1.2pt,
        fill=srsLight!80, text=srsDark,
        align=center, font=\small\sffamily,
        inner sep=8pt, minimum width=3.2cm, minimum height=1.4cm,
        drop shadow={opacity=0.2, shadow xshift=1.5pt, shadow yshift=-1.5pt}
    },
    systemNode/.style={
        rectangle, rounded corners=8pt,
        draw=srsPrimary, line width=1.8pt,
        fill=srsPrimary!12, text=srsDark,
        align=center, font=\normalsize\bfseries\sffamily,
        inner sep=12pt, minimum width=5.8cm, minimum height=3.2cm,
        drop shadow={opacity=0.25, shadow xshift=2pt, shadow yshift=-2pt}
    },
    externalNode/.style={
        rectangle, rounded corners=6pt,
        draw=srsAccent, line width=1.2pt,
        fill=srsAccent!10, text=srsDark,
        align=center, font=\small\sffamily,
        inner sep=8pt, minimum width=3.2cm, minimum height=1.4cm,
        drop shadow={opacity=0.2, shadow xshift=1.5pt, shadow yshift=-1.5pt}
    },
    flowArrow/.style={
        -{Stealth[scale=1.1]}, line width=1.2pt, draw=srsDark
    },
    biFlowArrow/.style={
        {Stealth[scale=1.1]}-{Stealth[scale=1.1]}, line width=1.2pt, draw=srsDark
    },
    labelNode/.style={
        font=\scriptsize\sffamily, text=srsDark!85, fill=white, inner sep=2pt, rounded corners=2pt
    }
]

    % Central System
    \node[systemNode] (plugin) at (0, 0) {
        \textbf{Midnight Compact Plugin}\\
        \vspace{4pt}
        \footnotesize\textmd{dev.verloren.midnight}\\
        \scriptsize\textnormal{IntelliJ Platform Language Extension}
    };

    % Left: Human Actors
    \node[actorNode] (dev) at (-6.2, 1.8) {
        \textbf{Smart Contract Dev}\\
        \scriptsize Authors .compact files,\\
        \scriptsize triggers refactorings \& builds
    };

    \node[actorNode] (auditor) at (-6.2, -1.8) {
        \textbf{ZK Security Auditor}\\
        \scriptsize Inspects circuits, sealed state,\\
        \scriptsize and witness disclosure flows
    };

    % Top: IDE Platform
    \node[externalNode] (idea) at (0, 4.0) {
        \textbf{JetBrains IntelliJ Platform}\\
        \scriptsize Editor, PSI Infrastructure, Indexing,\\
        \scriptsize Tool Windows, Gutter \& Annotations
    };

    % Right: External Compilers & Toolchains
    \node[externalNode] (cli) at (6.2, 1.8) {
        \textbf{Compact Toolchains}\\
        \scriptsize \texttt{\textasciitilde/.compact/versions/<ver>/compact}\\
        \scriptsize Isolated binaries (v0.31.1--0.34.0)
    };

    \node[externalNode] (wsl) at (6.2, -1.8) {
        \textbf{WSL2 / Linux Bridge}\\
        \scriptsize \texttt{/usr/bin/compact}\\
        \scriptsize Linux subsystem cross-execution
    };

    % Bottom: Midnight Network Ecosystem
    \node[externalNode] (midnight) at (0, -4.0) {
        \textbf{Midnight Blockchain Ecosystem}\\
        \scriptsize ZK-SNARK Verifier, Ledger Node,\\
        \scriptsize Proof Generation Tooling
    };

    % Connectors & Interactions
    \draw[flowArrow] (dev) -- node[labelNode, above, sloped] {Edits code / Ctrl+Space} (plugin.150);
    \draw[flowArrow] (auditor) -- node[labelNode, below, sloped] {Reviews ZK inspections} (plugin.210);

    \draw[biFlowArrow] (idea) -- node[labelNode, right] {PSI Events / Action System} (plugin);

    \draw[biFlowArrow] (plugin.30) -- node[labelNode, above, sloped] {CLI Invocation / Output} (cli);
    \draw[biFlowArrow] (plugin.330) -- node[labelNode, below, sloped] {WSL Process Execution} (wsl);

    \draw[flowArrow] (plugin) -- node[labelNode, right] {Compiled Artifacts / Proof Specs} (midnight);

\end{tikzpicture}

    \caption{System Context Diagram (C4 Model Level 1) for the Midnight Compact Language Plugin.}
    \label{fig:context_diagram}
\end{figure}

\subsection{References and Standards}
\begin{itemize}
    \item ISO/IEC/IEEE 29148:2018 -- Systems and software engineering -- Life cycle processes -- Requirements engineering~\cite{iso29148}.
    \item IEEE Std 830-1998 -- IEEE Recommended Practice for Software Requirements Specifications~\cite{ieee830}.
    \item RFC 2119 -- Key words for use in RFCs to Indicate Requirement Levels~\cite{rfc2119}.
    \item Midnight Network Architecture Whitepaper~\cite{midnightWhitepaper}.
    \item Compact Language Reference and Specification~\cite{compactSpec}.
    \item JetBrains IntelliJ Platform SDK Developer Guide~\cite{intellijSDK}.
    \item Top-Down Operator Precedence (Pratt Parsing)~\cite{pratt1973}.
    \item Semantic Versioning 2.0.0 Specification~\cite{semver2}.
\end{itemize}
